% Part: normal-modal-logic
% Chapter: completeness
% Section: lindenbaums-lemma

\documentclass[../../../include/open-logic-section]{subfiles}

\begin{document}

\olfileid{nml}{com}{lin}

\olsection{లిండెన్‌బామ్ ఉపసిద్ధాంతం}

లిండెన్‌బామ్ ఉపసిద్ధాంతం ప్రకారం, ప్రతి $\Sigma$-అవిరుద్ధ
\tetoken{సూత్రాల}{!!{formula}s} సమితి కనీసం ఒక
\emph{సంపూర్ణ} $\Sigma$-అవిరుద్ధ సమితిలో ఇమిడి ఉంటుంది.
కానానికల్ నమూనా నిర్మాణంలో, ప్రతి సంపూర్ణ $\Sigma$-అవిరుద్ధ
సమితి~$\Delta$కు ఆ నమూనాలో ఒక లోకం ఉంటుందని చూపుతాం;
ఆ లోకంలో సరిగ్గా $\Delta$లోని \tetoken{సూత్రాలు}{!!{formula}s}
మాత్రమే సత్యం. కనుక ప్రతి $\Sigma$-అవిరుద్ధ సమితి
కానానికల్ నమూనాలోని ఏదో ఒక లోకంలో సత్యమని
లిండెన్‌బామ్ ఉపసిద్ధాంతం హామీ ఇస్తుంది.

\begin{thm}[లిండెన్‌బామ్ ఉపసిద్ధాంతం]\ollabel{thm:lindenbaum}
  $\Gamma$ $\Sigma$-అవిరుద్ధమైతే, $\Gamma$ను విస్తరించే
  సంపూర్ణ $\Sigma$-అవిరుద్ధ సమితి $\Delta$ ఒకటి ఉంటుంది.
\end{thm}

\begin{proof}
$!A_0$, $!A_1$, \dots{} అని భాషలోని అన్ని సూత్రాల
సమగ్ర జాబితాను తీసుకుందాం (పునరావృత్తులు ఉండవచ్చు). ఉదాహరణకు,
ముందుగా $\Obj p_0$ను జాబితాలో చేర్చి, ప్రతి $n \ge 1$
దశలో పొడవు గరిష్ఠంగా~$n$ ఉండి, $\Obj p_0$,
\dots,~$\Obj p_n$లలోని చరరాశులను మాత్రమే ఉపయోగించే
పరిమిత సంఖ్యలోని సూత్రాలన్నిటినీ చేర్చవచ్చు.
\sourcecorrection{OLTENMLCOMLIN-001}{స్థిర మూలంలో ప్రతి $n$
  దశకు పొడవు సరిగ్గా $n$ ఉన్న సూత్రాలనే చేర్చమని ఉంది. అలా
  చేస్తే తక్కువ పొడవు గల, ఎక్కువ సూచికగల చరరాశిని వాడే
  సూత్రాలు జాబితా నుంచి తప్పిపోతాయి. పొడవు గరిష్ఠంగా $n$
  అని తీసుకుంటే ప్రతి దశ పరిమితంగానే ఉండి, జాబితా సమగ్రమవుతుంది.}
ఇప్పుడు $n$పై ఆగమనంతో \tetoken{సూత్రాల}{!!{formula}s}
సమితులు $\Delta_n$ను నిర్వచించి, తరువాత $\Delta = \bigcup_n
\Delta_n$ అని ఉంచుదాం. మొదట $\Delta_0 = \Gamma$ అని
ఉంచుదాం. $\Delta_n$ నిర్వచించబడిందని అనుకుని,
$\Delta_{n+1}$ను ఇలా నిర్వచిద్దాం:
\[
\Delta_{n+1} =
\begin{cases}
  \Delta_n \cup \{!A_n\}, & \text{$\Delta_n \cup \{ !A_n\}$
    $\Sigma$-అవిరుద్ధమైతే;} \\
  \Delta_n \cup \{ \lnot !A_n\}, & \text{లేకపోతే.}
\end{cases}
\]
ఇప్పుడు $\Delta = \bigcup_{n=0}^\infty \Delta_n$ అనుకుందాం.

ఈ నిర్వచనం నిజంగానే కావలసిన ధర్మాలుగల $\Delta$ సమితిని
ఇస్తుందని, అంటే $\Gamma \subseteq \Delta$ మరియు $\Delta$
సంపూర్ణ $\Sigma$-అవిరుద్ధమని, చూపాలి.

నిర్మాణం వల్ల $\Delta_0 \subseteq \Delta$, అలాగే
$\Delta_0 = \Gamma$; కాబట్టి $\Gamma \subseteq \Delta$
స్పష్టమే. నిజానికి $\Delta$ అన్ని $\Delta_n$ల సమ్మేళనం
కనుక ప్రతి~$n$కు $\Delta_n \subseteq \Delta$. (నిర్మాణంలోని
ప్రతి దశలో ఇప్పటికే నిర్మించిన సమితికి
\tetoken{ఒక సూత్రాన్ని}{!!a{formula}} చేర్చుతాం; కనుక
$\Delta_n \subseteq \Delta_{n+1}$. $\subseteq$ సంక్రమణీయత
వల్ల $n \le m$ అయినప్పుడల్లా $\Delta_n \subseteq
\Delta_{m}$.) ప్రతి దశలో $!A_n$ లేదా $\lnot !A_n$ను
చేర్చుతాం; ప్రతి \tetoken{సూత్రం}{!!{formula}} అన్ని~$!A_n$ల
జాబితాలో (కనీసం ఒకసారైనా) వస్తుంది. కనుక ప్రతి $!A$కు
$!A \in \Delta$ లేదా $\lnot !A \in \Delta$; నిర్వచనం
ప్రకారం $\Delta$ సంపూర్ణం.

చివరగా $\Delta$ $\Sigma$-అవిరుద్ధమని చూపాలి. దీనికోసం
(a) $\Delta$ $\Sigma$-విరుద్ధమైతే, ఏదో ఒక $\Delta_n$
$\Sigma$-విరుద్ధమవుతుందని, (b) అన్ని $\Delta_n$లు
$\Sigma$-అవిరుద్ధమని చూపుతాం.

$\Delta$ $\Sigma$-విరుద్ధమని అనుకుందాం. అప్పుడు $\Delta
\Proves[\Sigma] \lfalse$; అంటే, $!A_1$, \dots,~$!A_k \in
\Delta$ ఉండి, $\Sigma \Proves !A_1 \lif (!A_2 \lif \cdots (!A_k
\lif \lfalse)\dots)$. $\Delta = \bigcup_{n=0}^\infty \Delta_n$
కనుక, ప్రతి $!A_i \in \Delta_{n_i}$ అయ్యే ఏదో ఒక~$n_i$ ఉంటుంది.
ఈ సూచికలలో అతి పెద్దదాన్ని $n$ అనుకుందాం. $n_i \le n$ కనుక
$\Delta_{n_i} \subseteq \Delta_n$. అందువల్ల ఆ $!A_i$లన్నీ
ఒకే~$\Delta_n$లో ఉంటాయి. అప్పుడు $\Delta_n
\Proves[\Sigma] \lfalse$; అంటే $\Delta_n$ $\Sigma$-విరుద్ధం.

ప్రతి $\Delta_n$ $\Sigma$-అవిరుద్ధమని చూపడానికి $n$పై
సరళ ఆగమనం వాడుదాం. $\Delta_0 = \Gamma$; $\Gamma$
$\Sigma$-అవిరుద్ధమని మన పరికల్పన. కనుక $n = 0$కు వాదన
నిజం. ఇప్పుడు $n$కు నిజమని, అంటే $\Delta_n$
$\Sigma$-అవిరుద్ధమని అనుకుందాం. $\Delta_n \cup \{!A_n\}$
$\Sigma$-అవిరుద్ధమైతే $\Delta_{n+1}$ అదే; లేకపోతే
$\Delta_n \cup \{\lnot!A_n\}$ అవుతుంది. మొదటి సందర్భంలో
$\Delta_{n+1}$ $\Sigma$-అవిరుద్ధమని స్పష్టమే. అయితే
\olref[prf][con]{prop:consistencyfacts}\olref[prf][con]{prop:consistencyfacts-c}
ప్రకారం $\Delta_n \cup \{!A_n\}$ లేదా $\Delta_n \cup
\{\lnot!A_n\}$ అవిరుద్ధమే. కనుక రెండవ సందర్భంలోనూ
$\Delta_{n+1}$ అవిరుద్ధం.
\end{proof}

\begin{cor}\ollabel{cor:provability-characterization}
  $\Gamma \Proves[\Sigma] !A$ అయితే, అప్పుడు మాత్రమే
  $\Gamma$ను విస్తరించే ప్రతి సంపూర్ణ $\Sigma$-అవిరుద్ధ
  సమితి $\Delta$లో $!A \in \Delta$ (ఇది $\Gamma = \emptyset$
  అయినప్పుడూ వర్తిస్తుంది; అప్పుడు మోడల్ వ్యవస్థ $\Sigma$కు
  మరొక లక్షణీకరణ లభిస్తుంది).
\end{cor}

\begin{proof}
  $\Gamma \Proves[\Sigma] !A$ అనుకుని, $\Gamma$ను విస్తరించే
  ఏదైనా సంపూర్ణ $\Sigma$-అవిరుద్ధ సమితి $\Delta$ను తీసుకుందాం.
  $!A \notin \Delta$ అయితే గరిష్ఠత వల్ల $\lnot!A \in \Delta$.
  అప్పుడు ఏకదిశత్వం వల్ల $\Delta \Proves[\Sigma] !A$;
  స్వప్రతిఫలనత్వం వల్ల $\Delta \Proves[\Sigma] \lnot!A$.
  కాబట్టి $\Delta$ విరుద్ధమవుతుంది. విపరీతంగా,
  $\Gamma \Proves/[\Sigma] !A$ అయితే $\Gamma \cup \{
  \lnot!A\}$ $\Sigma$-అవిరుద్ధం. లిండెన్‌బామ్
  ఉపసిద్ధాంతం ప్రకారం $\Gamma \cup \{ \lnot!A \}$ను
  విస్తరించే సంపూర్ణ అవిరుద్ధ సమితి $\Delta$ ఒకటి ఉంటుంది.
  అవిరుధ్యం వల్ల $!A \notin \Delta$.
\end{proof}

\end{document}
